\documentclass[12pt]{article}
\usepackage[utf8]{inputenc}
\usepackage[T1]{fontenc}
\usepackage{hyperref}
\title{T\"ubatulabal Language Resources for K-12 Curriculum Development}
\author{Compiled by Claude Sonnet 5 (claude-sonnet-5) for Andr\'as Kornai}
\date{}
\begin{document}
\maketitle

\section{Introduction}
T\"ubatulabal (ISO 639-3: tub; also Pahka'anil) is a Uto-Aztecan language of Kern River Valley, California. Wikipedia lists it as extinct (July 30, 2008, on the death of Jim Andreas), while UNESCO classifies it as critically endangered. Revitalization work by community members is reported.

\section{Primary Educational Resources}
No school materials were found. A 2004 linguist-list job posting described community work on a writing system and on dictionaries, grammars and language lessons.

\section{Linguistic and Reference Materials}
\begin{itemize}
\item Overview: \url{https://en.wikipedia.org/wiki/T%C3%BCbatulabal_language}
\item Voegelin, C. F., Working Dictionary of T\"ubatulabal (IJAL, 1958), described at: \url{https://wiki.ubc.ca/Documentation:RelLex/Working_Dictionary_of_T%C3%BCbatulabal}
\item Voegelin, C. F., T\"ubatulabal grammar, Univ. of California Publications in American Archaeology and Ethnology 34 (1935), pp. 55--190; foundational description. Catalogue record: \url{https://search.ub.tum.de/vufind/Record/DE-604.BV021077820}
\item Marean et al. (2021), Pahka'anil (T\"ubatulabal): Two Texts, IJAL; no URL seen.
\item Berkeley Language and Culture Atlas: \url{https://cla.berkeley.edu/languages/tubatulabal.html}
\item Job posting on community documentation work: \url{https://listserv.linguistlist.org/pipermail/linguist/2004-December/025749.html}
\end{itemize}

\section{Religious and Community Resources}
None were found.

\section{Digital and Audio Resources}
Endangered Languages Project records: \url{https://endangeredlanguages.com/node/109518} and \url{https://endangeredlanguages.com/node/105067}. Smithsonian archive entry: \url{https://si.edu/object/archives/components/sova-naa-1976-95-ref14231}.

\section{Curriculum Development Context}
The language is documented mainly in mid-twentieth-century descriptions and archival recordings. No current revitalization programme details were found; a heritage curriculum would build on the Voegelin materials and any community work.

\end{document}